% Supplementary research design; not executed results.
% Current convention: Y(m)=1 means unresolved; S(m)=1-Y(m) means success.
\paragraph{Outcome convention.} As in the current main paper, $Y(m)=1$ denotes terminal failure under action $m$, and $S(m)=1-Y(m)$ denotes successful resolution. Proposed intervention utility is $Q_m(h_t)=\mathbb{E}[S(m)-\lambda C(m)-\mu D(m)\mid h_t]$, with costs $C$, adverse effects $D$, and nonnegative weights $\lambda,\mu$. This is an unexecuted design, not an estimated result.
\section{Toward an Intervention-Aware PreAct}
\label{sec:toward-an-intervention-aware-preact}
\subsection{Actionable Event Labels}
An eventual-failure label is insufficient to identify a prevention window. A next-stage dataset should annotate three separate times: the first prefix with observable warning evidence, the first consequential agent error, and the earliest point after which recovery is infeasible under the task budget. These boundaries need executable checks where available, with adjudicated review for ambiguous cases. A discrete-time hazard model can then estimate failure risk over several future-step horizons while accounting for runs that terminate before a horizon. In particular, an alert raised after the consequential error should not count as prevention even when the terminal task has not yet failed.

\subsection{Matched-Prefix Branches}
The central unresolved estimand is the value of each monitor action at the \emph{same} live prefix. For a task and fixed environment checkpoint, record the proposed next tool action and branch the run into three initial policies: continue without monitor input, request one low-cost verification step, or restart with the same base agent. Each branch receives the same total token and tool budget, including the monitor's overhead. Repeated stochastic seeds expose variance. If external tools cannot be replayed consistently, the branch must be marked non-comparable rather than treated as a controlled counterfactual. Following prior executed repair studies~\cite{sentinel,dover}, a subsequent online trial should randomize eligible prefixes to monitor policies and analyze outcomes by task cluster.

The primary outcome is the difference in task-resolution probability at matched budget. Secondary outcomes include added tokens, tool calls, wall time, the fraction of originally successful runs disrupted by intervention, and the fraction of originally unsuccessful runs recovered. This design addresses the disruption--recovery tradeoff observed by Vasudev et al.~\cite{interventionparadox}. A failure classifier may be used to propose branch candidates, but selection must be judged by the downstream task outcome, not by classifier accuracy alone.

\subsection{Separating Model And Scaffold Shift}
The follow-up should independently vary model and scaffold. A factorial design trains on one combination and evaluates a new model with the original scaffold, the original model with a new scaffold, and their joint change. Raw structural, normalized, semantic and event-based monitors should share task splits and budgets. A target calibration set must remain separate from target testing. This distinguishes harness shortcuts from model change, which our current pooled transfer test cannot do. Organic failures should be the primary cohort; injected faults provide a separate robustness arm. Since matched-checkpoint repair and recalibration already exist~\cite{sentinel,dover}, the intended advance is a prospectively tested, cost-constrained policy under these controlled shifts, not branching alone.

\subsection{Calibrated Decision And Abstention}
For each feasible action $m$, estimate $\Delta_m(h_t)=Q_m(h_t)-Q_{\mathrm{continue}}(h_t)$ from matched branch data. The monitor should intervene only when a lower confidence bound for net benefit exceeds zero under a prespecified false-disruption and cost budget. Otherwise it continues. New agent systems require a separate calibration check on held-out tasks; a confidence interval or out-of-support detector should trigger abstention when the monitor lacks evidence. Our observed transfer failure makes this a necessary part of the evaluation, not merely an optional deployment refinement. The study should compare this value-based policy with fixed failure-risk thresholds, reactive recovery after explicit errors, and single-action restart or verification policies.

